\section{Actuarial impact}
\label{sec:actuarial}

\subsection{H5: propagation into life expectancy and annuity factors}

\inputtable{tab-h5-actuarial-main}

Table~\ref{tab:h5-actuarial-main} gives, by arm and regime, nominal-95\%
coverage and mean error for $e_0$, $e_{65}$ and $\annuity$ at 2\%, from latent
rate paths (Section~\ref{sec:metrics-actuarial}); Table~S6 carries
the full fifty-cell fragment. Conformal cells emit no latent paths, so every
count here is over the twenty distributional arms. CBD is fit on ages 55--99:
it contributes no $e_0$, and its $e_{65}$ and $\annuity$ are exact on a
truncated table whose fit never saw the younger ages. Excluding error rows
(Table~S11), no row carries its regime's full complement---520 stable, 40
shift, 22 placebo (population, sex, origin) units: the fullest stable rows
reach 442, the sparse VAR native 297 and bootstrap 238, the GP sixteen of
twenty stable populations, the sparse VAR's native placebo cell 0 of 22.
Addendum~3~\S8 restricts \Hfive\ to ex-post coverage of these derived
intervals, with attenuation---integration averaging a rate failure away
\citep{cairns2011mortality}---as the named competitor.

\Hfive\ registered that interval miscalibration propagates: coverage of 95\%
intervals for $e_{65}$ and $\annuity$ departs from nominal by more than the
corresponding stable-regime departure. As the universal ordering it states, it fails: the shift departure exceeds the stable one for 14 of 20 arms
on $e_{65}$ and 12 of 20 on $\annuity$---a majority, not a universal, by the
standard that declines \Htwo's registered form at 27 of 50
(Section~\ref{sec:results}). What survives is this audit's most consequential
actuarial finding: miscalibration reaches the reserve in every regime, the
derived intervals already broken in the stable control. Against nominal 0.95,
control coverage of realised $e_{65}$/$\annuity$ runs: SVD Lee--Carter native
0.256/0.253 on 442 units, LSTM ensemble 0.182/0.157, LSTM dropout 0.205/0.191,
Poisson-LC 0.434/0.430, neural-LC ensemble 0.468/0.459, CNN-LC dropout
0.559/0.548, APC/RH 0.654/0.643. Only two arms come within 0.05 of nominal on
$\annuity$---the sparse VAR native at 0.980 on 297 units and CBD's bootstrap at
0.934 on its 55--99 support---and thirteen of twenty sit below 0.700. An
interval containing the realised value in a quarter of origins is not a usable
reserving input, in 1990--2014 as much as across 2020.

The registered comparison orders two-sided departures
$|\text{coverage}-0.95|$, where the result is partial. The six exceptions on
$e_{65}$---both SVD Lee--Carter and both LSTM arms, CNN-LC dropout, Poisson-LC
bootstrap---are joined on $\annuity$ by CNN-LC's ensemble and Poisson-LC's
native law. Mostly these are not intervals the break spared but the
largest stable departures moved \emph{towards} nominal: SVD Lee--Carter's
bootstrap 0.645 to 0.400 on $\annuity$. The rule does not separate
cleanly---Poisson-LC native and the neural-LC ensemble start further from
nominal on $e_{65}$ (0.516 and 0.482) than CNN-LC dropout (0.391) and still
move away---and \Htwo's rate-scale sorting cannot operate, since nineteen of
twenty arms already under-cover $\annuity$ in the control and the one that does
not, the sparse VAR native, is the arm the break damages most, 0.030 to 0.364.
A floor operates instead: the eight span 0.157 to 0.558 in the control, and an
interval already missing four cells in five has little room to miss more.

The attenuation alternative fails in both regimes. Against the marginal rate
coverage of Table~S3, all eighteen full-age distributional arms cover realised
$\annuity$ less often than a single log rate; CBD's two truncated arms are
excluded and run the other way. The sides differ in horizon---marginal
coverage pools $h=1\ldots5$, the functional uses the first test year alone---so
read the sizes with that mismatch in view; the 18-of-18 direction does not turn
on it, the annuity being scored at its most favourable horizon. The gap runs
from 0.008 (sparse VAR native) to 0.481 (LSTM ensemble), averaging 0.252 in the
control and 0.289 under COVID. Integration does not average the failure away;
it concentrates it, in quiet periods too. Some cells look calibrated on rates and
not on liabilities: in the control CNN-LC dropout covers 0.954 of single log
rates against 0.548 of annuity factors, the NB head's
ensemble 0.988 against 0.814.

The mechanism is \Hthree's on the age axis: a period functional integrates a
whole age range, so a common-factor miss that leaves single-age intervals
intact still displaces it---one displaced functional, one uncovered cell.
Priced from the same marginal baseline on the same eighteen arms, the age axis
is dearer: marginal-to-joint-path gaps average 0.153 in the control and 0.208
under shift (Tables~S3 and S4) against the 0.252 and 0.289 above. The 0.127 and
0.168 quoted for \Hthree\ (Section~\ref{sec:results}) average over all fifty
cells, thirty conformal with no derived quantity, so are not like-for-like. The
horizon gap is also the weaker quantity, most of it arithmetic: independence at
nominal coverage over five horizons gives $1-0.95^5=0.176$, above either grid
mean. The comparator is that independence benchmark, not each model's own implied
path law: Addendum~3~\S8's re-specified comparator, a cell's model-implied
joint coverage, is not computable from the emitted columns---recorded as a
limitation. The annuity gap admits no arithmetic discount:
$\annuity$ is one event on one test year.

Two families qualify the mechanism. Only CBD has
derived coverage above its own rate-scale coverage: in the control its native
law covers 0.930 of realised $e_{65}$ against a marginal 0.901 and its
bootstrap 0.950 and 0.934 against a marginal 0.921, the native arm's 0.889 on
$\annuity$ the one figure of four falling short. CBD is also the only family
fit and integrated on ages 55--99; family and truncated support are not
separable here. The sparse VAR is the one full-age arm calibrated when nothing
breaks (0.983 and 0.980) and the arm the break damages most (0.621 and 0.586,
on 29 admissible shift units); its scales disagree---rates move \emph{towards}
nominal, 0.988 to 0.937, while the liability falls to 0.586---so one scale
alone would call it improved or destroyed at the analyst's choice. On
$e_0$, which \Hfive\ does not register, Poisson-LC covers 0.744 and APC/RH
0.808 in the control, above their rate-level marginals of 0.682 and 0.731,
while SVD Lee--Carter reaches 0.369 and the LSTM ensemble 0.267; that column is
descriptive.

\subsection{Reading a coverage shortfall in balance-sheet terms}

The interval on $\annuity$ is an interval on the best-estimate liability of a
unit annuitant. In Solvency~II terms an internal model covering the realised
value in materially fewer than 95\% of population--sex cells has a
miscalibrated longevity distribution---too narrow for most arms here,
mislocated for the few that overstate---and Table~\ref{tab:h5-actuarial-main}
puts every distributional arm there under shift, all but two in the control:
under COVID the best is 0.850 (CBD, truncated table), no full-age arm above
0.750.

In the control the mean error on $\annuity$ is negative for 19 of 20 arms---the
central factor below the realised one, understating the liability---the NB-head
ensemble excepted at $+0.02$ annuity units; the understatement reaches $-0.63$
for the LSTM ensemble and $-0.58$ for SVD Lee--Carter, whose median $e_{65}$
path sits a mean 0.84 years below the realised value in the control, 0.57 below
at the break. Across the pandemic the signs split, twelve positive to
eight negative: both SVD Lee--Carter, both CNN-LC and both LSTM arms still sit
below the realised $\annuity$ ($-0.24$ to $-0.45$), CBD's two marginally below
($-0.03$ and $-0.04$), while APC/RH, neural-LC, the sparse VAR native and the
NB-head ensemble sit above ($+0.24$ to $+0.39$), the sparse VAR bootstrap by
$+0.14$. This is not the single-signed effect a uniform mortality
margin would absorb, and for the arms already understating it is no pandemic
effect at all: the excess mortality of 2020 moved the realised factor towards
their too-low estimates, masking a pre-existing bias.

A 97.5\% quantile below the realised $\annuity$ understates implied capital by
the shortfall, reported relative to the central $\annuity$. These breach
frequencies replace the registered 99.5\% longevity quantile, which the runner
did not emit (\S8). The understating arms breach at many times the nominal
2.5\%, more often in the control than under COVID: there the realised
$\annuity$ exceeds the 97.5\% quantile in 0.804 of the LSTM ensemble's 434
cells, 0.779 of LSTM dropout's, 0.701 of SVD Lee--Carter's 442, 0.529 of
Poisson-LC's and 0.441 of CNN-LC dropout's, at mean shortfalls of 4.2\%, 4.0\%,
4.2\%, 2.0\% and 2.4\%; the same five breach in 0.605, 0.474, 0.425, 0.150 and
0.325 of shift cells, at 3.4\%, 3.8\%, 4.3\%, 1.5\% and 2.5\%. Capital at the
97.5\% quantile is thus understated by two to four per cent of the
best-estimate liability in a large minority of origins---for four of these
five, a control majority---against a nominal one in forty. The overstating arms
fail on the other side: under COVID the NB-head ensemble and its dropout arm
never breach above, but the realised value falls below their 2.5\% quantile in
0.550 and 0.250 of cells (mean excess 2.0\% and 1.6\%), and below the neural-LC
ensemble's in 0.625 (2.3\%); in the control the three are milder and two-sided,
upper breaches of 0.081, 0.106 and 0.292 against lower breaches of 0.104, 0.090
and 0.249---reserves too high rather than capital too low. CBD breaches above
in 0.050--0.075 of shift and 0.066--0.109 of stable cells at a mean shortfall
at or below 1\%.

\subsection{Twin crises in economic terms}

The placebo rows give the same functionals for 1914--1922, descriptively: with
the control in place, a shortfall the quiet decades already show cannot be
credited to either break. For the full-age classical cores under native laws,
realised $e_{65}$ falls inside the band in 0.425--0.500 of
population--sex cells under COVID and 0.455--0.636 across 1914--1922,
$\annuity$ in 0.450--0.475 and 0.455--0.636---against 0.256--0.654 and
0.253--0.643 in the control, on 442 units against 40 and 22. The shortfall is
material at both breaks, of the same order at each, and for two of three cores
no larger than the control's: SVD Lee--Carter's annuity coverage is
\emph{lowest} in the control, 0.253, highest at the placebo, 0.636, COVID
between at 0.450; Poisson-LC runs the same way at 0.430, 0.475 and 0.591.
APC/RH is the exception, worse at both breaks, its $e_{65}$ and $\annuity$
falling from 0.654 and 0.643 to 0.500 and 0.475 under COVID and to 0.455 at the
placebo. CBD's native law covers 0.875 on $e_{65}$ under COVID and 1.000 on
both functionals at the placebo against 0.889 and 0.930 in the control, its
age-55 support above the excess mortality of 1918. The $e_0$ column shows
the timing caveat---0.909 for Lee--Carter at the placebo against 0.500 under
shift and 0.369 in the control---because 1914, the placebo's first test year,
predates the influenza wave while 2020 sits inside the COVID shock. Outside the
sparse VAR's placebo bootstrap row---1.000 on both functionals, on the nine of
twenty-two units that survived---no full-age family's derived intervals
are near nominal at either break, nor in the control: across three regimes the pandemic is not distinctive in economic kind, nor, for the worst
families, in degree.
